\Askhsh{13767}{2}
\begin{ylh}{1}
    \item 3.1. Είδη και στοιχεία τριγώνων
    \item 3.2. 1ο Κριτήριο ισότητας τριγώνων
    \item 4.6. Άθροισμα γωνιών τριγώνου
    \item 5.4. Ρόμβος.
\end{ylh}
% ===================================================================
%  Αρχείο: 13767.tex (Fragment)
%  Αυτόματη μετατροπή μέσω doc2latex.py
% ===================================================================

ΘΕΜΑ 2

Σε ευθεία ε θεωρούμε τα διαδοχικά σημεία Α, Β και Γ έτσι ώστε ΑΒ$\  = \ $ΒΓ. Στη συνέχεια, σχεδιάζουμε τα ισόπλευρα τρίγωνα ΑΒΔ και ΒΓΕ προς το ίδιο ημιεπίπεδο ως προς την ευθεία ε όπως φαίνεται στο παρακάτω σχήμα.

\begin{alist}
\item Να υπολογίσετε τη γωνία $\text{Δ}\widehat{\text{Β}}\text{Ε}$.

\monades{7}

\item Να αποδείξετε ότι το τρίγωνο ΒΔΕ είναι ισόπλευρο.

\monades{10}

\item Να αποδείξετε ότι το τετράπλευρο ΑΔΕΒ είναι ρόμβος.

\monades{8}

\includesvg[width=4.88868in,height=2.04724in]{/home/spyros/Μαθηματικά/Trapeza_Thematwn/A_Lykeiou/Geometry/extracted/13767/media/image2.svg}
\end{alist}
